\documentclass[11pt,a4paper]{article}
\usepackage{amsmath}
\usepackage{amssymb}
\usepackage{amsfonts}
\usepackage{graphicx}
\usepackage{hyperref}
\usepackage{geometry}
\usepackage{caption}
\usepackage{subcaption}
\usepackage{booktabs}
\usepackage{xcolor}
\usepackage{titlesec}
\usepackage{float}

\geometry{margin=1in}

\title{Nichols Radial Injection Model (RIM) VIII\\The Bullet Cluster and Geometric Torsion}
\author{Lawrence William Nichols}
\date{March 29, 2026}

\begin{document}

\maketitle

\begin{abstract}
The Radial Injection Model (RIM) eliminates the need for exotic dark-matter particles. All apparent dark-sector phenomenology arises from purely mechanical interactions between a dynamic 3D hypersurface and a static 4D bulk. Radial ERB mass flux $\Phi_m$ drives viscous surface friction, reproducing the baryonic Tully–Fisher relation $v^4 \propto M_{\rm baryon}$ and excess gravitational lensing with zero free parameters. Global geometric torsion inherited from the 5D line element
\[
ds^2 = R^2 \left( d\psi^2 + \sin^2 \psi \,(d\theta^2 + \sin^2 \theta \,d\phi^2) \right) + (dt + \omega \,dw)^2,
\]
with torsion constant $\omega = 0.004$, produces a closed Round Corridor ($L = 4.65\,\rm Tly$, $v_{\rm spin} \approx 2047\,\rm km\,s^{-1}$). This torsional spin supplies the collisionless behaviour observed in the Bullet Cluster while simultaneously inducing Relativistic Baryonic Latency: injected baryonic mass is gravitationally active on the hypersurface today, yet its electromagnetic signal remains delayed by the full torsional cycle. The Little Red Dots observed by JWST are the first photons from these ancient injections finally arriving. The hypersphere itself thus supplies the entire dark sector through a single, calibrated 5D mechanism.
\end{abstract}

\section{Dark Matter as Hypersurface Friction and Torsional Spin}

The Radial Injection Model (RIM) derives all apparent dark-sector effects from purely mechanical interactions between a dynamic 3D hypersurface and a static 4D bulk.

\begin{enumerate}
    \item Viscous surface friction as the 3-sphere expands radially via ERB mass flux $\Phi_m$.
    \item Global torsional spin inherited from the 5D metric, combined with temporal decoupling of injected matter.
\end{enumerate}

The 5D line element of the manifold is
\[
ds^2 = R^2 \left( d\psi^2 + \sin^2 \psi \,(d\theta^2 + \sin^2 \theta \,d\phi^2) \right) + (dt + \omega \,dw)^2,
\]
where the torsion constant $\omega = 0.004$ couples temporal drift to the extra dimension $w$. This produces a closed ``Round Corridor'' along the longitudinal orbit $L = 4.65\,\rm Tly$ with transverse radius $r_{\rm trans} = 5 \times 10^9\,\rm ly$.

The effective tangential spin speed from one full 4-rotation cycle is
\[
v_{\rm spin} = c \cdot \frac{2\pi r_{\rm trans}}{L} \approx 2047\,\rm km\,s^{-1}.
\]

\subsection{Friction on Galactic Scales}

On galactic scales, the hypersurface slides outward through the static 4D bulk, producing viscous drag on baryonic matter:
\[
F_{\rm fric} \propto -\Phi_m \cdot v_{\rm surface},
\]
where $v_{\rm surface}$ is the local expansion velocity of the 3-sphere. For a star in circular orbit the force balance becomes
\[
\frac{v^2}{r} = \frac{GM_{\rm baryon}(r)}{r^2} + k\Phi_m v,
\]
with $k$ fixed by $\omega = 0.004$. In the outer, friction-dominated regime this yields the baryonic \textbf{Tully--Fisher relation}
\[
v^4 \propto M_{\rm baryon}
\]
without free parameters. The same friction also enhances light deflection, explaining excess gravitational lensing around baryonic matter. Both effects arise mechanically from the expanding hypersurface interacting with the static bulk.

\subsection{Collisionless Behaviour in Cluster Mergers}

In high-velocity cluster collisions such as the Bullet Cluster (1E 0657-56), the global torsional spin and associated temporal decoupling dominate. The observed velocity of the subcluster’s lensing mass centroid ($\sim 2600\,\rm km\,s^{-1}$) is close to the geometric spin speed $v_{\rm spin} \approx 2047\,\rm km\,s^{-1}$, while the intracluster gas shock propagates faster ($\sim 4500$--$4700\,\rm km\,s^{-1}$).

Because the torsional component resides in the 5D bulk geometry, it is unaffected by electromagnetic ram pressure. This produces the observed separation:
\begin{itemize}
    \item \textbf{Baryonic Gas}: Interacts via ram pressure, slows, and piles up at the centre (seen in X-rays).
    \item \textbf{Lensing Mass}: Carried collisionlessly by geometric torsion, continuing on a straight trajectory.
    \item \textbf{Stars and Galaxies}: Remain coupled to the hypersurface and follow the lensing signal.
\end{itemize}
No separate collisionless dark-matter fluid is needed; the hypersphere itself supplies the dark sector.

\subsection{Primordial Decoupling and the First Injection}

The present 3-sphere results from radial ERB mass flux $\Phi_m$ from the static 4D bulk. Any transition from higher-dimensional geometry to the hypersurface involves an energy release as matter couples to the lower-dimensional manifold. Gamma-ray bursts (GRBs) are a plausible signature of this process: we observe the intense electromagnetic flash and afterglow (the visible ``effects''), but the injected baryonic mass/energy itself decouples temporally from the local hypersurface coordinate.

In contrast to Segal’s chronometric cosmology—where redshift arises geometrically from light propagation across a static curved manifold without energy loss—the RIM inverts the picture. The injected baryons become ``overly fast.'' Over the closed longitudinal orbit $L = 4.65\,\rm Tly$ their invariant mass remains gravitationally active on the 3-sphere via energy density $E = mc^2$, while their electromagnetic history is delayed by the full torsional cycle.

The temporal lag is
\[
\Delta t_{\rm lag} = \oint_L \frac{dw}{c(1-\omega)},
\]
where $\omega = 0.004$. This lag reaches trillions of years.

Even in the Big Bang era (the earliest phase of injection), speeds were so extreme that the electromagnetic signals from that primordial matter have still not caught up with the hypersurface’s local time coordinate. What is conventionally called ``dark matter'' is therefore \textbf{Relativistic Baryonic Latency}: mass that is physically present and gravitationally active today, yet whose light originates from a time coordinate deep in the primordial epoch.

During Bullet Cluster-like mergers, this temporally decoupled component (the accumulated result of past injections) sails through the collision zone carried by geometric torsion at speeds near $v_{\rm spin}$. It remains unaffected by ram pressure while the braked baryonic gas piles up centrally. The observed separation of lensing mass from hot gas is a direct geometric consequence of the hypersphere’s 4-rotation, surface friction, and the energy release accompanying injections from the 4D bulk. The hypersphere itself, seeded by ERB injections whose GRB signatures we occasionally detect, supplies the entire dark sector.

\section{Video Catalog}

A video overview of the full RIM series is available on the author’s Nichols Radial Injection Model (RIM) YouTube channel.

\section*{References}

\begin{enumerate}
    \item L.~W. Nichols. \textit{Nichols Radial Injection Model (RIM) and Radial ERB Inflows: A Mechanism for Progenitor-less Astrophysical Events, the Hubble Pulse, and 4D Substrate Interactions.} [Self-published], January 21, 2026. \url{https://rxiverse.org/abs/2602.0036}
    \item L.~W. Nichols. \textit{Nichols Radial Injection Model (RIM) 11-Year Solar Cycle Pulses, GRB Milestones, Historical Alignments, and the 16.6 Gyr Hypersphere Universe.} [Self-published], February 2026. \url{https://rxiverse.org/abs/2602.0038}
    \item L.~W. Nichols. \textit{Nichols Radial Injection Model (RIM) III: The 3,000-Year Universal Odometer, the 1054 CE Primary Strike, and the Decaying 11-Year Cycle.} [Self-published], 15 February 2026. \url{https://rxiverse.org/abs/2602.0043}
    \item L.~W. Nichols. \textit{Nichols Radial Injection Model (RIM) IV: A 5D Mass-Regulated Manifold Solution to Cosmic Expansion.} [Self-published], 28 February 2026. \url{https://rxiverse.org/abs/2603.0001}
    \item L.~W. Nichols. \textit{Nichols Radial Injection Model (RIM) V: The Empirical Seal of the 11.07-Year Pulse and Universal Resonance in Stellar Cycles.} [Self-published], 4 March 2026. \url{https://rxiverse.org/abs/2603.0003}
    \item L.~W. Nichols. \textit{Nichols Radial Injection Model (RIM) VI: The Birth of the Universe.} [Self-published], 10 March 2026. \url{https://rxiverse.org/abs/2603.0023}
    \item L.~W. Nichols. \textit{Nichols Radial Injection Model (RIM) VII: The $\Lambda$CDM Inconsistencies.} [Self-published], 24 March 2026. \url{https://rxiverse.org/abs/2603.0085}
    \item A. Einstein \& N. Rosen. (1935). The Particle Problem in the General Theory of Relativity. \textit{Physical Review}, 48(1), 73.
\end{enumerate}

\section*{Historical Precedent and Author Bios}

\textbf{Lawrence William Nichols} is a researcher in astrophysics and the architect of the Radial Injection Model (RIM). His work focuses on 5D manifold mechanics and the elimination of non-baryonic dark matter through geometric torsion.

\textbf{Irving Ezra Segal} (1918--1998) was a mathematician known for his Chronometric Cosmology. While his model relied on static geometry, RIM acknowledges his foundational insight that universal observation is a product of manifold curvature rather than local expansion alone.

\section{Conclusion}

The Radial Injection Model provides a unified mechanical explanation for dark-sector phenomena through a single engine: radial ERB injections $\Phi_m$ driving volumetric growth and energetic transients, surface friction shaping galactic dynamics and lensing, and torsional spin plus relativistic baryonic latency producing collisionless behaviour in mergers. This framework accounts for flat rotation curves, the Tully--Fisher relation, excess lensing, and the Bullet Cluster’s mass-gas separation—all consistent with the calibrated parameters ($\omega = 0.004$, $L = 4.65\,\rm Tly$, $\xi = 20$).

The hypersphere acts as the structural 4D substrate supplying both dark-matter-like and dark-energy-like effects without exotic particles. Even the extreme conditions of the early universe fit naturally: matter injected at such high speeds that its electromagnetic signals have not yet caught up, leaving a gravitationally active but temporally lagged population that we observe today only through its gravitational footprint.

\section{Engineering Studies and Visual Data}

\newpage

\section{Engineering Studies and Visual Data}
This section includes relevant visual diagrams.
\includepdf[pages=-]{Geometric_Dark_Matter.pdf}

\end{document}